\chapter{Feynman diagram interpretation of bar-cobar duality}
\label{ch:feynman}

\begin{remark}[Logarithmic Feynman--bar comparison]
\label{rem:feynman-bar-bridge}
\index{Feynman diagram!bar-cobar bridge|textbf}
Let $\cA$ be a chiral algebra on a smooth curve~$X$ with OPE
$a(z)\,b(w) \sim \sum_n C^{ab}_n(w)\,(z-w)^{-n-1}$.
The bar complex records the same local collision data as the
configuration-space Feynman expansion.  The edge form is fixed in
residue normalization:
\[
\eta_{ij}=d\log(z_i-z_j),\qquad
\operatorname{PRes}_{D_{ij}}(\eta_{ij})=1,\qquad
\frac{1}{2\pi i}\int_{\gamma_{ij}}\eta_{ij}=1,
\]
where $D_{ij}=\{z_i=z_j\}$ and $\gamma_{ij}$ is a small positive normal
circle to~$D_{ij}$.
\begin{center}
\renewcommand{\arraystretch}{1.25}
\begin{tabular}{lcl}
\textbf{Configuration-space datum} && \textbf{Bar datum $\barBch(\cA)$} \\
\hline
diagonal singularity $z_i=z_j$
 && logarithmic edge form $\eta_{ij}$ \\
vertex factor $C^{ab}_n$
 && OPE residue $\Res_{z_i = z_j}[a(z_i)\,b(z_j)\,(z_i-z_j)^n]$ \\
short-distance boundary
 && FM boundary stratum $\partial_S\overline{C}_k(X)$, $|S|\ge 2$ \\
global compactification
 && compactification $C_k(X) \hookrightarrow \overline{C}_k(X)$
\end{tabular}
\end{center}
Theorem~\ref{thm:bar-cobar-isomorphism-main}
identifies the bar differential with residues of the logarithmic forms
$\eta_{ij} = d\log(z_i-z_j)$ along collision divisors; the bar
complex carries logarithmic edge kernels and OPE vertex
residues. The Fulton--MacPherson compactification replaces
short-distance collisions by normal-crossing boundary strata.
Theorem~\ref{thm:loop-genus-formula} gives the relation between
graph loop number $h^1(\Gamma)$ and ribbon-surface genus, which
also depends on the number of boundary faces.
\end{remark}

\index{Feynman diagram|textbf}

\section{Feynman diagrams in chiral field theory}

\begin{definition}[Chiral graph weights]
\label{def:chiral-graph-weights}
\index{propagator!Feynman}
\index{Feynman diagram!and bar complex}
A \emph{chiral graph weight} on a curve~$X$ is built from a chiral
algebra~$\cA$, residue-normalized edge forms
\(\eta_{ij}=d\log(z_i-z_j)\) on a disk chart, and the OPE residue maps
of~\(\cA\).  On a higher-genus curve~\(\Sigma\), the local edge form is
replaced by \(d_{z_i}\log E_\Sigma(z_i,z_j)\), with \(E_\Sigma\) the
prime form; its diagonal residue is again~\(1\).

For a stable graph~\(\Gamma\) with ordered external legs, define
\[
\Omega_\Gamma(\phi_1,\ldots,\phi_k)
=\frac{\epsilon(\Gamma)}{|\Aut(\Gamma)|}
\left(\bigwedge_{e=(i,j)\in E(\Gamma)}\eta_{ij}\right)
\otimes
\left(\bigotimes_{v\in V(\Gamma)}C_v\right)(\phi_1,\ldots,\phi_k),
\]
on the corresponding Fulton--MacPherson stratum.  The sign
\(\epsilon(\Gamma)\) is the orientation sign of the stratum and
\(|\Aut(\Gamma)|^{-1}\) is the stable-graph symmetry factor.  When the
form degree matches the dimension of the chosen chain~\(\sigma_\Gamma\),
the scalar graph contribution is
\[
A_\Gamma(\phi_1,\ldots,\phi_k)
=\int_{\sigma_\Gamma}\Omega_\Gamma(\phi_1,\ldots,\phi_k).
\]
\end{definition}

\subsection{Tree vs.\ loop decomposition}

\begin{definition}[Loop number]
A Feynman diagram $\Gamma$ with $V$ vertices, $E$ edges, and $C$ connected
components has loop number
\[L(\Gamma) = E - V + C,\]
the first Betti number $b_1(\Gamma)$ of the graph.
\end{definition}

\begin{remark}[Configuration space interpretation]
The number \(L(\Gamma)\) counts independent cycles in the graph.  It
does not by itself count configuration variables: before quotienting by
automorphisms, \(V\) vertices move in \(C_V(X)\), and a boundary stratum
imposes collision relations according to its stable tree or stable
graph.  Surface genus is obtained only by adding ribbon or modular
surface data, as in Theorem~\ref{thm:loop-genus-formula}.
\end{remark}

\section{Bar-cobar residue pairing}

\begin{definition}[Bar-cobar distributional pairing]
\label{def:physical-pairing}
Let \(D\subset \overline{C}_n(X)\) be a union of collision divisors and
let \(\iota\colon C_n(X)\hookrightarrow \overline{C}_n(X)\) be the open
configuration-space inclusion.  For
\[
\omega_{\mathrm{bar}}\in
\Omega^\bullet_{\log}(\overline{C}_n(X),\cA^{\boxtimes n}),
\qquad
K_{\mathrm{cobar}}\in
H^\bullet_D(\overline{C}_n(X),\cA^{!\boxtimes n}),
\]
the Verdier residue pairing is
\[
\langle \omega_{\mathrm{bar}},K_{\mathrm{cobar}}\rangle_D
=\operatorname{Tr}_{\cA,\cA^!}
\operatorname{PRes}_{D}
\bigl(\omega_{\mathrm{bar}}\wedge K_{\mathrm{cobar}}\bigr).
\]

\begin{center}
\begin{tabular}{lll}
\toprule
\textbf{Boundary datum} & \textbf{Bar complex} & \textbf{Cobar complex} \\
\midrule
External leg & Marked point & Dual marked point \\
Internal edge & Logarithmic form $\eta_{ij}$ & Local cohomology class on $D_{ij}$ \\
Vertex & OPE residue & Dual coproduct \\
Boundary stratum & Poincar\'e residue & Support condition \\
Symmetry factor & $|\Aut(\Gamma)|^{-1}$ & $|\Aut(\Gamma)|^{-1}$ \\
\bottomrule
\end{tabular}
\end{center}
\end{definition}

\begin{example}[Heisenberg edge normalization]
For the rank-one Heisenberg current \(J(z)J(w)\sim (z-w)^{-2}\), the
bar edge
\[
\omega = J(z_1) \otimes J(z_2) \otimes \eta_{12}
\in \Omega^1(\overline{C}_2(X), \cA^{\boxtimes 2})
\]
carries the logarithmic pole \(\eta_{12}=d\log(z_1-z_2)\).  The
Verdier-dual edge is the local cohomology class
$K \in H^1_{D_{12}}(\overline{C}_2(X), \cA^{!\boxtimes 2})$
supported on the diagonal \(D_{12}=\{z_1=z_2\}\).  With the local
generator normalized by \(K(\eta_{12})=1\),
\[
\langle \omega,K\rangle_{D_{12}}=1.
\]
The OPE coefficient of the double pole and the logarithmic edge
residue are separate normalizations: \(J(z)J(w)\sim (z-w)^{-2}\)
fixes the vertex coefficient, while
\(\operatorname{PRes}_{D_{12}}\eta_{12}=1\) fixes the propagator edge.
\end{example}

\section{Higher operations from boundary strata}

\subsection{Virasoro fourth-pole normalization}

\begin{example}[Virasoro fourth-pole residue]\label{ex:virasoro-fourth-pole}
For the Virasoro algebra,
\[T(z) = \sum_n L_n z^{-n-2}, \quad [L_m, L_n] = (m-n)L_{m+n} +
\frac{c}{12}m(m^2-1)\delta_{m+n,0}\]
and
\[T(z)T(w) = \frac{c/2}{(z-w)^4} + \frac{2T(w)}{(z-w)^2} +
\frac{\partial T(w)}{z-w}+\mathrm{reg}.
\]
The ordinary one-variable residue of the fourth-order pole is
\[
\operatorname{Res}_{u=0}
\left(\frac{c/2}{u^4}\,T(w+u)\,du\right)
=\frac{c}{2}\cdot\frac{1}{3!}\partial^3T(w)
=\frac{c}{12}\partial^3T(w).
\]
Equivalently, in the divided-power \(\lambda\)-bracket convention,
\[
\{T_\lambda T\}
=\partial T+2\lambda T+\frac{c}{12}\lambda^3,
\qquad
T_{(3)}T=c/2.
\]
Thus the two scalar normalizations used in the Feynman lane are
distinct: the Virasoro curvature parameter is
\(\kappaChHodge(\mathrm{Vir}_c)=c/2\), while the divided-power
\(\lambda^3\)-coefficient is \(c/12\).
\end{example}

\subsection{Higher operations and factorization}

\begin{theorem}[\texorpdfstring{$A_\infty$}{A-infinity} constraint formula]
\label{thm:ainfty-constraint-formula}
The $A_\infty$ operations satisfy the recursive constraint:
\[m_1 \circ m_n = -\sum_{\substack{r+s+t=n \\ (r,s,t) \neq (0,n,0)}} (-1)^{rs+t} m_{r+1+t}(\mathrm{id}^{\otimes r} \otimes m_s \otimes \mathrm{id}^{\otimes t})\]

This constrains $m_1(m_n)$ (not $m_n$ itself) in terms of lower operations,
with the same boundary-incidence pattern as forest subtractions after a
choice of comparison map to renormalized amplitudes.
\end{theorem}

\begin{proof}
This is a rewriting of the $A_\infty$ relations (Keller convention):
\[\sum_{r+s+t=n} (-1)^{rs+t} m_{r+1+t}(\mathrm{id}^{\otimes r} \otimes m_s \otimes
\mathrm{id}^{\otimes t}) = 0.\]
Isolating the $r = t = 0$, $s = n$ term, which is $m_1(m_n)$, gives the result.
To \emph{define} $m_n$ recursively, for example by homotopy transfer,
requires a contracting homotopy $h$ with $h m_1 + m_1 h = \mathrm{id} - \iota\pi$;
the displayed constraint is a \emph{consistency relation} among the operations,
not a recursive definition.
\end{proof}

\section{Graph complexes and Kontsevich formality}

\subsection{The graph complex}

\begin{definition}[Kontsevich graph complex]
The graph complex $\text{GC}_n$ consists of:
\begin{itemize}
\item Generators: Isomorphism classes of graphs with $n$ labeled external legs
and unlabeled internal vertices
\item Differential: Sum over edge contractions
\item Grading: Loop number $L(\Gamma)$
\end{itemize}
\end{definition}

\begin{definition}[Decorated graph-to-bar map]
\label{def:bar-graph-map}
Let \(\mathrm{GC}_X(\cA)\) be the span of finite graphs whose external
legs are labelled by fields of~\(\cA\), whose internal vertices are
labelled by OPE residue operations, and whose edges carry the Koszul
orientation sign.  The map
\[
\Phi_{\Gamma\to\barB}\colon \mathrm{GC}_X(\cA)\longrightarrow
\Omega^\bullet_{\log}(\overline{C}_\bullet(X),\cA^{\boxtimes \bullet})
\]
sends a decorated graph~\(\Gamma\) to the stable-graph form
\(\Omega_\Gamma\) of Definition~\ref{def:chiral-graph-weights}.  For an
edge \(e=(i,j)\) contracted to the diagonal \(D_{ij}\), the edge
contraction differential is carried to the Poincar\'e residue:
\[
\Phi_{\Gamma/e\to\barB}(\Gamma/e)
=\operatorname{PRes}_{D_{ij}}\Phi_{\Gamma\to\barB}(\Gamma),
\qquad
\operatorname{PRes}_{D_{ij}}\eta_{ij}=1.
\]
The remaining terms in the bar differential are the vertex OPE residue
and the boundary-stratum splitting maps.
\end{definition}

\subsection{Kontsevich's formality and chiral algebras}

\begin{theorem}[Kontsevich formality]
\label{thm:kontsevich-formality-feynman}
For a smooth manifold $M$, the $L_\infty$ algebra of polyvector fields
$\mathcal{T}_{\text{poly}}(M)$ is formal (\cite{Kon03}):
\[\mathcal{T}_{\text{poly}}(M) \simeq_{L_\infty} H^*(\mathcal{T}_{\text{poly}}(M))\]
\end{theorem}

\begin{remark}[Kontsevich weights and chiral logarithmic weights]
\label{rem:kontsevich-graphs}
\label{sec:connections-other-feynman}
Kontsevich's formality weights integrate angle forms on compactified
configuration spaces~\cite{Kon99}.  The chiral graph map
\(\Phi_{\Gamma\to\barB}\) uses the logarithmic forms
\(\eta_{ij}=d\log(z_i-z_j)\).  The common structure is the
edge-form assignment plus Stokes' theorem on configuration-space
boundaries; the chiral target is the bar-cobar complex of the chiral
algebra rather than the polyvector-field Hochschild complex.
\end{remark}

\begin{definition}[Disk-local perturbative/FM discrepancy]
\label{def:disk-local-perturbative-fm}
Let $U \subset X$ be a holomorphic disk chart on a curve, and suppose a
holomorphic-topological perturbative field theory on~$U$ produces the
same local chiral algebra model \(\cA|_U\).  Let
\(\beta_2^{\mathrm{pert}}\) and \(\beta_3^{\mathrm{pert}}\) denote the
regularized BRST-anomaly brackets of Gaiotto--Kulp--Wu~\cite{GKW24}
on this local model, and let \(m_2^{\mathrm{bar}}\) and
\(m_3^{\mathrm{bar}}\) be the bar operations cut out by the
Fulton--MacPherson collision strata.  Define
\[
\Delta_2=\beta_2^{\mathrm{pert}}-m_2^{\mathrm{bar}},
\qquad
\Delta_3=\beta_3^{\mathrm{pert}}-m_3^{\mathrm{bar}},
\]
where \(m_3^{\mathrm{bar}}\) is the oriented sum of the residues on the
three codimension-\(1\) boundary divisors of \(\overline{C}_3(U)\).
The disk-local comparison is the vanishing of
\((\Delta_2,\Delta_3)\).  Its quadratic obstruction is the difference
between the Wess--Zumino relation for the perturbative brackets and the
Stokes/\(A_\infty\) relation for the corresponding bar operations.
\end{definition}

\begin{remark}[FM extension class]\label{rem:disk-local-fm-extension}
The disk-local forms define an FM logarithmic comparison precisely when
their local representatives extend to
\(\Omega^\bullet_{\log}(\overline{C}_n(U))\) with residues on all
collision strata.  The obstruction is the relative logarithmic
extension class
\[
\mathrm{ob}_{\mathrm{FM}}(\omega_U)\in
H^1\!\left(\overline{C}_n(U),
\Omega^\bullet_{\log}(\partial\overline{C}_n(U))
/\Omega^\bullet_{\log,\mathrm{ext}}(\omega_U)\right).
\]
When this class vanishes, Stokes regularity on FM compactifications
\textup{(}Proposition~\textup{\ref{V1-prop:stokes-regularity-FM}}\textup{)}
identifies the quadratic boundary relation with the bar
\(A_\infty\)-relation.
\end{remark}

\begin{proposition}[Binary-ternary reduction for the disk-local discrepancy]
\label{prop:disk-local-binary-ternary-reduction}
Fix a holomorphic disk chart $U \subset X$, and interpret the
perturbative BRST-anomaly brackets and the local bar operations on the
same local chain model underlying~$\cA|_U$.
Assume:
\begin{enumerate}[label=\textup{(\roman*)}]
\item the binary perturbative bracket agrees with the local bar
 residue/OPE operation on~$C_2(U)$; and
\item the ternary perturbative bracket agrees with the local
 $m_3$-operation, equivalently with the sum of the three
 codimension-$1$ residues on~$\overline{C}_3(U)$.
\end{enumerate}
Then the quadratic/Wess--Zumino identities on the perturbative side
agree automatically with the quadratic $A_\infty$/Stokes relation on
the bar side. In particular, there is no third independent local
coefficient calculation beyond the binary and ternary comparisons.
\end{proposition}

\begin{proof}
On the perturbative side, Gaiotto--Kulp--Wu~\cite{GKW24} prove that the
regularized BRST-anomaly brackets satisfy the required quadratic
axioms/Wess--Zumino consistency relations.

On the bar side, the corresponding degree-$3$ consistency relation is
already built into the local $A_\infty$ structure: by
Theorem~\ref{thm:ainfty-constraint-formula} specialized to $n=3$, the
bar operations satisfy the quadratic relation expressing $m_1m_3$ in
terms of composites of lower operations; geometrically this is the
Stokes relation among the codimension-$1$ boundary strata of the
compactified three-point configuration space.

Once the binary and ternary operations are
identified on the common local chain model, the two quadratic packages
compare the same formal relation; no third
independent local coefficient problem remains.
\end{proof}

\begin{remark}[Application to the disk-local FM comparison]
Applied to Definition~\ref{def:disk-local-perturbative-fm},
Proposition~\ref{prop:disk-local-binary-ternary-reduction} shows that
the only independent local input is the binary comparison on $C_2(U)$
and the ternary comparison on $\overline{C}_3(U)$.
\end{remark}

\begin{corollary}[Anomaly-free genus-\texorpdfstring{$0$}{0} binary vanishing]
\label{cor:disk-local-ternary-on-brstbar-locus}
On the anomaly-free genus-$0$ BRST/bar locus of
Theorem~\ref{thm:brst-bar-genus0}, the binary discrepancy
\(\Delta_2\) of Definition~\ref{def:disk-local-perturbative-fm}
vanishes.  Hence the disk-local comparison is equivalent on this locus
to the vanishing of the ternary discrepancy \(\Delta_3\) on
\(\overline{C}_3\).
\end{corollary}

\begin{proof}
Theorem~\ref{thm:brst-bar-genus0} identifies the genus-$0$ BRST complex
with the genus-$0$ bar complex by a chain map
\(\Phi\) to \(\barB^{\mathrm{ch}}_0(\cA_{\mathrm{tot}})\).
Remark~\ref{rem:brst-bar-explicit} makes the low-degree content of that
map concrete: physical vertex operators map to bar generators, and the
ghost bilinear maps to the logarithmic propagator
\(\eta_{12}=d\log(z_1-z_2)\).
Thus the propagator/residue data underlying the binary discrepancy is
identified on this genus-$0$ locus.

By Proposition~\ref{prop:disk-local-binary-ternary-reduction}, once the
binary discrepancy vanishes the local comparison is controlled by the
ternary discrepancy.
\end{proof}

\begin{proposition}[Two-channel reduction after compactifying the
 ternary discrepancy]
\label{prop:compactified-ternary-two-channel}
Fix a holomorphic disk chart $U \subset X$, and suppose the
perturbative ternary bracket and the local bar $m_3$-operation are
represented by logarithmic $1$-forms on
$\overline{M}_{0,4} \cong \bP^1$ whose only poles lie at the three
boundary points corresponding to the channels
$(12|3)$, $(23|1)$, and $(13|2)$.
Then comparison of the two residues in any two of those channels
implies comparison in the third channel as well.
Equivalently, after compactification the ternary discrepancy has two
independent boundary-channel checks.
\end{proposition}

\begin{proof}
Let $\omega_{\mathrm{pert}}$ and $\omega_{\mathrm{bar}}$ denote the two
logarithmic $1$-forms on $\overline{M}_{0,4} \cong \bP^1$, and set
\[
\delta = \omega_{\mathrm{pert}} - \omega_{\mathrm{bar}}
\in H^0\!\left(\bP^1,\,
\Omega^1_{\bP^1}\bigl(\log\{0,1,\infty\}\bigr)\right).
\]
If the residues of $\delta$ agree in two channels, then two of the
three residues of $\delta$ vanish. By the residue theorem on~$\bP^1$,
the sum of all three residues is zero, so the third residue vanishes as
well. Therefore $\delta$ has zero residue at every boundary point, so
it is a holomorphic $1$-form on~$\bP^1$. But
$H^0(\bP^1,\Omega^1_{\bP^1}) = 0$, hence $\delta = 0$.
The third boundary residue is forced by the first two.
\end{proof}

\begin{remark}[Application to the disk-local ternary discrepancy]
Proposition~\ref{prop:compactified-ternary-two-channel} applies in
particular to the compactified ternary discrepancy of
Definition~\ref{def:disk-local-perturbative-fm}. Once the two local
ternary forms are represented on $\overline{M}_{0,4}$, vanishing of
two boundary residues forces vanishing of the third.
\end{remark}

\begin{corollary}[Genus-\texorpdfstring{$0$}{0} post-compactification ternary reduction]
\label{cor:genus0-compactified-ternary-two-channel}
If, on the anomaly-free genus-$0$ BRST/bar locus of
Theorem~\ref{thm:brst-bar-genus0}, the ternary discrepancy is
represented by logarithmic \(1\)-forms on
$\overline{M}_{0,4} \cong \bP^1$, then \(\Delta_3=0\) is equivalent
to vanishing of any two of its three boundary residues.
\end{corollary}

\begin{proof}
Corollary~\ref{cor:disk-local-ternary-on-brstbar-locus} reduces the
disk-local discrepancy on this genus-$0$ locus to \(\Delta_3\).
Proposition~\ref{prop:compactified-ternary-two-channel} identifies the
vanishing of \(\Delta_3\) with any two boundary-channel residues.
\end{proof}

\begin{corollary}[Standard-chart form of the genus-\texorpdfstring{$0$}{0}
 ternary discrepancy]
\label{cor:genus0-standard-chart-two-residues}
Choose the standard coordinate on
$\overline{M}_{0,4} \cong \bP^1$ with boundary points
$0,1,\infty$. On the anomaly-free genus-$0$ BRST/bar locus of
Theorem~\ref{thm:brst-bar-genus0}, the compactified ternary discrepancy
vanishes exactly when its residues at $0$ and~$1$ vanish; the residue at
$\infty$ is then forced.
\end{corollary}

\begin{proof}
By Corollary~\ref{cor:genus0-compactified-ternary-two-channel}, any two
boundary-channel checks suffice. On the standard
$0,1,\infty$ compactification of $\overline{M}_{0,4}$, choose
the channels at $0$ and~$1$ as the independent pair. Then
Proposition~\ref{prop:compactified-ternary-two-channel} forces the
residue at~$\infty$.
\end{proof}

\begin{proposition}[Standard logarithmic basis on
 \texorpdfstring{$\overline{M}_{0,4}$}{M_0,4}]
\label{prop:m04-standard-log-basis}
Let $t$ be the standard coordinate on
$\overline{M}_{0,4} \cong \bP^1$ with boundary points
$0,1,\infty$. Then every logarithmic $1$-form on
$\bP^1$ with poles only at $\{0,1,\infty\}$ admits a unique
expression
\[
\omega = A\, d\log t + B\, d\log(1-t)
\]
for scalars $A,B \in \bC$.
Equivalently, the residues at $0$ and $1$ determine the form
uniquely, and the residue at $\infty$ equals $-A-B$.
\end{proposition}

\begin{proof}
Write $\omega = f(t)\,dt$ with $f(t)$ a rational function on~$\bP^1$
having at most simple poles at $0$, $1$, and $\infty$.
Thus
\[
f(t)=\frac{a}{t}+\frac{b}{t-1}+c
\]
for some scalars $a,b,c$. The condition that $\omega$ have at most a
simple pole at $\infty$ forces $c=0$: in the coordinate $u=1/t$ one
has
\[
dt = -u^{-2}\,du,\qquad
\left(\frac{a}{t}+\frac{b}{t-1}+c\right)dt
= \left(\frac{a}{u^{-1}}+\frac{b}{u^{-1}-1}+c\right)(-u^{-2}\,du),
\]
and a nonzero constant term $c$ would contribute a double pole at
$u=0$.
Hence
\[
\omega = a\,\frac{dt}{t}+b\,\frac{dt}{t-1}
= a\, d\log t - b\, d\log(1-t).
\]
Setting $A=a$ and $B=-b$ gives existence.

For uniqueness, if
$A\, d\log t + B\, d\log(1-t)=0$, then the residues at
$t=0$ and $t=1$ are respectively $A$ and $B$, so $A=B=0$.
The residue at $\infty$ is minus the sum of the finite residues, hence
$-A-B$.
\end{proof}

\begin{corollary}[Two-coefficient form of the genus-\texorpdfstring{$0$}{0}
 compactified discrepancy]
\label{cor:genus0-two-coefficient-discrepancy}
On the anomaly-free genus-$0$ BRST/bar locus of
Theorem~\ref{thm:brst-bar-genus0}, the compactified ternary discrepancy
is exactly the vanishing of two scalar coefficients. More precisely, if
\[
\delta = \omega_{\mathrm{pert}} - \omega_{\mathrm{bar}}
\in H^0\!\left(\bP^1,\,
\Omega^1_{\bP^1}\bigl(\log\{0,1,\infty\}\bigr)\right),
\]
then uniquely
\[
\delta = A\, d\log t + B\, d\log(1-t),
\]
and \(\delta=0\) if and only if \(A=B=0\).
\end{corollary}

\begin{proof}
Corollary~\ref{cor:genus0-standard-chart-two-residues} identifies the
vanishing condition with the two residues at $0$ and $1$ in the
standard chart. Proposition~\ref{prop:m04-standard-log-basis}
identifies those residues with the coefficients of \(d\log t\) and
\(d\log(1-t)\).
\end{proof}

\begin{definition}[Standard genus-\texorpdfstring{$0$}{0} coefficient pair]
\label{def:genus0-coefficient-pair}
For
\[
\omega \in H^0\!\left(\bP^1,\,
\Omega^1_{\bP^1}\bigl(\log\{0,1,\infty\}\bigr)\right),
\]
write uniquely
\[
\omega=\mathsf{a}(\omega)\,d\log t+\mathsf{b}(\omega)\,d\log(1-t)
\]
as in Proposition~\ref{prop:m04-standard-log-basis}. For the
compactified ternary perturbative and bar forms, write
\[
\omega_{\mathrm{pert}}
= \mathsf{a}_{\mathrm{pert}}\,d\log t
 + \mathsf{b}_{\mathrm{pert}}\,d\log(1-t),
\qquad
\omega_{\mathrm{bar}}
= \mathsf{a}_{\mathrm{bar}}\,d\log t
 + \mathsf{b}_{\mathrm{bar}}\,d\log(1-t).
\]
Call the pair of equalities
\[
\mathsf{a}_{\mathrm{pert}}=\mathsf{a}_{\mathrm{bar}},
\qquad
\mathsf{b}_{\mathrm{pert}}=\mathsf{b}_{\mathrm{bar}}
\]
the \emph{standard genus-$0$ coefficient pair}.
\end{definition}

\begin{corollary}[Named coefficient form of the genus-\texorpdfstring{$0$}{0}
 compactified discrepancy]
\label{cor:genus0-named-coefficient-discrepancy}
On the anomaly-free genus-$0$ BRST/bar locus of
Theorem~\ref{thm:brst-bar-genus0}, the compactified ternary discrepancy
vanishes exactly when the standard genus-$0$ coefficient pair of
Definition~\ref{def:genus0-coefficient-pair} holds.
\end{corollary}

\begin{proof}
Corollary~\ref{cor:genus0-two-coefficient-discrepancy} identifies the
compactified genus-$0$ discrepancy with the vanishing of the two
coefficients of
\(\delta=\omega_{\mathrm{pert}}-\omega_{\mathrm{bar}}\).
This is equivalent to the two equalities
\[
\mathsf{a}_{\mathrm{pert}}=\mathsf{a}_{\mathrm{bar}},
\qquad
\mathsf{b}_{\mathrm{pert}}=\mathsf{b}_{\mathrm{bar}}.
\]
\end{proof}

\section{\texorpdfstring{The $m_k$ operations as transfer graph sums}{The m k operations as transfer graph sums}}
\label{sec:mk-feynman-complete}

\subsection{\texorpdfstring{Arity and transfer degree}{Arity and transfer degree}}

\begin{definition}[The complete \texorpdfstring{$m_k$}{m_k} transfer family]
\label{def:mk-family-feynman}
Let
\[
\bigl(H(\mathcal A),0\bigr)
\xrightarrow{\ \iota\ }
(\mathcal A,m_1)
\xrightarrow{\ \pi\ }
\bigl(H(\mathcal A),0\bigr),
\qquad
h m_1+m_1h=\id-\iota\pi,
\]
be a contraction of the underlying complex.  The transferred
\(A_\infty\) operations
\(\widetilde m_k\colon H(\mathcal A)^{\otimes k}\to H(\mathcal A)\)
extend to coderivations of \(\bar B(H(\mathcal A))\); the component
induced by \(\widetilde m_k\) maps
\(\bar B^n\to \bar B^{n-k+1}\) for \(n\geq k\).

\begin{center}
\begin{tabular}{|>{\raggedright\arraybackslash}p{1.2cm}|>{\raggedright\arraybackslash}p{4.2cm}|>{\raggedright\arraybackslash}p{5.6cm}|>{\raggedright\arraybackslash}p{1.8cm}|}
\hline
\textbf{$k$} & \textbf{Algebraic datum} & \textbf{Transfer graph} & \textbf{$h$-edges} \\
\hline
\(m_0\) & curvature & nullary insertion & \(0\) \\
\hline
\(m_1\) & differential & one external edge & \(0\) \\
\hline
\(m_2\) & binary product/OPE residue & one binary vertex & \(0\) \\
\hline
\(m_3\) & associator homotopy & two binary vertices joined by one \(h\) & \(1\) \\
\hline
\(m_4\) & pentagon operation & five planar binary trees & \(2\) \\
\hline
\(m_k\) & \(k\)-ary transfer operation & planar rooted binary trees with \(k\) leaves & \(k-2\) \\
\hline
\end{tabular}
\end{center}
The last column counts internal homotopy edges in a genus-\(0\)
transfer tree.  It is not the genus of a surface.
\end{definition}

\begin{theorem}[Graph loop number and ribbon genus]
\label{thm:loop-genus-formula}
Let \(\Gamma\) be a finite graph with \(V\) vertices, \(E\) internal
edges, and \(C\) connected components.  Its loop number is
\[
\ell(\Gamma)=h^1(\Gamma)=E-V+C.
\]
If \(\Gamma\) is connected and is equipped with an orientable ribbon
structure whose thickening has \(F\) boundary components and genus
\(g\), then
\[
2g=\ell(\Gamma)+1-F.
\]
Consequently \(g\leq \ell(\Gamma)/2\), with equality exactly in the
one-boundary case \(F=1\).  A genus-\(0\) homotopy-transfer tree has
\(\ell(\Gamma)=0\).
\end{theorem}

\begin{proof}
The cellular chain complex of a graph gives
\(h^1(\Gamma)=E-V+C\).  For a connected ribbon graph, the thickened
surface \(\Sigma\) deformation-retracts to~\(\Gamma\), so
\[
\chi(\Sigma)=V-E=1-\ell(\Gamma).
\]
The same surface has \(\chi(\Sigma)=2-2g-F\).  Equating the two
Euler characteristics gives \(2g=\ell(\Gamma)+1-F\).  Since
\(F\geq 1\), \(g\leq \ell/2\), and equality is equivalent to
\(F=1\).
\end{proof}

\subsection{\texorpdfstring{$m_2$ and $m_3$ as boundary residues}{m2 and m3 as boundary residues}}

\begin{example}[Binary product as OPE residue]
\label{ex:m2-tree-level}
The binary operation is the diagonal residue
\[
m_2(\phi_1,\phi_2)
=\operatorname{PRes}_{D_{12}}
\bigl(\phi_1(z_1)\phi_2(z_2)\eta_{12}\bigr).
\]
If
\[
\phi_1(z)\phi_2(w)
\sim \sum_{r\geq 0}C_r(\phi_1,\phi_2)(w)(z-w)^{-r-1},
\]
then the coefficient \(C_r\) is extracted by the ordinary residue
\[
C_r(\phi_1,\phi_2)(w)
=\operatorname{Res}_{z=w}
\bigl((z-w)^r\phi_1(z)\phi_2(w)\,dz\bigr).
\]
The logarithmic edge normalization
\(\operatorname{PRes}_{D_{12}}\eta_{12}=1\) is independent of the
OPE pole order.
\end{example}

\begin{example}[Ternary operation as a boundary sum]
\label{ex:m3-boundary}
On the genus-\(0\) four-point compactification
\(\overline{M}_{0,4}\cong\bP^1\), the ternary operation is an oriented
sum over the three boundary channels:
\[
m_3(\phi_1,\phi_2,\phi_3)
=\sum_{D\in\{(12|3),(23|1),(13|2)\}}
\epsilon_D\,\operatorname{PRes}_D
\bigl(\omega_D(\phi_1,\phi_2,\phi_3)\bigr).
\]
The signs \(\epsilon_D\) are the boundary-orientation signs of
\(\overline{M}_{0,4}\).  This is a genus-\(0\) boundary formula; no
surface genus is assigned without the ribbon or stable-graph data of
Theorem~\ref{thm:loop-genus-formula}.
\end{example}

\begin{computation}[Virasoro \texorpdfstring{$m_3$}{m3} coefficient class]
\label{comp:virasoro-m3}
Fix transfer data \((\pi,\iota,h)\).  In the Keller sign convention the
ternary transferred operation has the form
\[
\widetilde m_3
=\pi\Bigl[
m_2(hm_2\otimes\id)
-m_2(\id\otimes hm_2)
\Bigr]\iota^{\otimes 3},
\]
with the global suspension convention used to identify coderivations
with operations fixed as in Theorem~\ref{thm:ainfty-constraint-formula}.
For the Virasoro input \(T^{\otimes 3}\), define the
\(\partial^3T\)-coefficient by projection
\[
\mu_3(h,\pi,\iota)
=\operatorname{coeff}_{\partial^3T}
\bigl(\widetilde m_3(T,T,T)\bigr).
\]
Each fourth-pole contribution entering this coefficient uses the
normalization of Example~\ref{ex:virasoro-fourth-pole}: the OPE pole is
\(T_{(3)}T=c/2\), and its divided-power \(\lambda^3\)-coefficient is
\(c/12\). Hence
\[
\mu_3(h,\pi,\iota)
\in \frac{c}{12}\,\bC,
\]
the scalar determined by the two oriented transfer trees and the
chosen contraction \(h\).
\end{computation}

\subsection{\texorpdfstring{$m_4$ and stable graph genus}{m4 and stable graph genus}}

\begin{remark}[\texorpdfstring{$m_4$}{m4} transfer degree and stable genus]
\label{rem:m4-transfer-stable-genus}
The genus-\(0\) transfer operation \(m_4\) is a sum over the five
planar rooted binary trees with four leaves.  Each such tree has two
internal homotopy edges and graph loop number \(0\).  A stable graph
with \(h^1(\Gamma)=1\) or \(2\) belongs instead to the modular
stable-graph expansion, where the total curve genus is
\[
g(\Gamma)=\sum_{v\in V(\Gamma)}g_v+h^1(\Gamma).
\]
Thus transfer depth, graph loop number, and curve genus are three
separate gradings.
\end{remark}

\begin{computation}[Virasoro \texorpdfstring{$m_4$}{m4} transfer cochain]
\label{comp:virasoro-m4}
\index{Virasoro algebra!$A_\infty$ operations!$m_4$}
The homotopy transfer theorem (Theorem~\ref{V1-thm:htt}) applied to
$\B(\mathrm{Vir}_c) \to H^*(\B(\mathrm{Vir}_c))$
produces operations $\tilde{m}_k = \pi \circ T_k \circ \iota^{\otimes k}$,
where $\pi$ is the projection to cohomology, $\iota$ the inclusion,
and $h$ the contracting homotopy (propagator).
For $m_4$, the recursive formula
(Appendix~\ref{V1-app:homotopy-transfer}) gives:
\begin{multline}\label{eq:m4-transfer}
\tilde{m}_4 = \pi \circ \Bigl[
m_2(hm_2(hm_2 \otimes \id)\otimes \id) +
m_2(hm_2(\id \otimes hm_2)\otimes \id) \\
{} + m_2(hm_2 \otimes hm_2) +
m_2(\id \otimes hm_2(hm_2 \otimes \id)) \\
{} + m_2(\id \otimes hm_2(\id \otimes hm_2))
\Bigr] \circ \iota^{\otimes 4}.
\end{multline}

\emph{The five trees.}
The five terms correspond to the vertices of the Stasheff
associahedron $K_4$ (pentagon), i.e., the $C_3 = 5$ planar binary
trees with $4$~leaves:
\begin{center}
\renewcommand{\arraystretch}{1.3}
\begin{tabular}{cll}
\toprule
$\Gamma$ & Parenthesization & Internal structure \\
\midrule
$\Gamma_1$ & $((T \cdot T) \cdot T) \cdot T$ & left comb \\
$\Gamma_2$ & $(T \cdot (T \cdot T)) \cdot T$ & left-right-left \\
$\Gamma_3$ & $(T \cdot T) \cdot (T \cdot T)$ & balanced \\
$\Gamma_4$ & $T \cdot ((T \cdot T) \cdot T)$ & right-left-right \\
$\Gamma_5$ & $T \cdot (T \cdot (T \cdot T))$ & right comb \\
\bottomrule
\end{tabular}
\end{center}
Each tree $\Gamma_i$ has $2$ internal edges (propagator insertions~$h$)
and $3$ internal vertices (binary products~$m_2$).

\emph{Conformal weight support.}
The input $T^{\otimes 4}$ has total conformal weight $4 \times 2 = 8$.
With the standard residue grading, two homotopy edges land in the
weight-\(6\) span
\begin{equation}\label{eq:m4-weight-analysis}
\widetilde m_4(T^{\otimes 4}) \;\in\; \mathbb{C}\cdot\partial^4 T \;\oplus\;
\mathbb{C}\cdot \partial^2({:}TT{:}) \;\oplus\;
\mathbb{C}\cdot {:}\partial T \partial T{:} \;\oplus\;
\mathbb{C}\cdot {:}T\partial^2 T{:} \;\oplus\;
\mathbb{C}\cdot {:}TTT{:},
\end{equation}
a \(5\)-dimensional space of weight-\(6\) fields. The central
charge~\(c\) enters only through the fourth-order pole
$T(z)T(w) \sim (c/2)(z-w)^{-4} + \cdots$,
so a double central contraction contributes a factor \((c/2)^2\)
before the divided-power and projection normalizations are applied.

\emph{The \(\partial^4T\)-coefficient.}
Define the coefficient
\begin{equation}\label{eq:m4-leading}
\mu_4^{\partial^4T}(h,\pi,\iota)
=\operatorname{coeff}_{\partial^4T}
\bigl(\widetilde m_4(T,T,T,T)\bigr).
\end{equation}
It is the finite signed tree sum
\[
\mu_4^{\partial^4T}(h,\pi,\iota)
=\sum_{\Gamma\in\mathrm{PBT}_4}
\epsilon_\Gamma\,
\operatorname{coeff}_{\partial^4T}
\left(\pi\,\mu_\Gamma(h,m_2)\,\iota^{\otimes 4}(T^{\otimes4})\right),
\]
where \(\mathrm{PBT}_4\) is the set of planar binary trees above and
\(\mu_\Gamma(h,m_2)\) is the composite dictated by~\(\Gamma\).  Planar
rooted trees have trivial planar automorphism group; if the trees are
quotiented to non-planar stable graphs the corresponding term is
weighted by \(1/|\Aut(\Gamma)|\).  Every fourth-pole residue uses the
normalizations \(T_{(3)}T=c/2\) and \(T_{(3)}T/3!=c/12\) from
Example~\ref{ex:virasoro-fourth-pole}; the scalar coefficient is the
displayed function \(\mu_4^{\partial^4T}(h,\pi,\iota)\) of the chosen
contraction and projection.

\emph{Pentagon identity.}
The $A_\infty$ relation at degree~$4$ reads
(Appendix~\ref{V1-app:homotopy-transfer},
equation~\eqref{V1-eq:ainfty-relation}):
\begin{equation}\label{eq:pentagon-ainfty}
\sum_{\substack{r+s+t = 4 \\ s \geq 2}} (-1)^{rs+t}\,
\tilde{m}_{r+1+t}\bigl(\id^{\otimes r} \otimes \tilde{m}_s
\otimes \id^{\otimes t}\bigr) = 0.
\end{equation}
The five terms are:
\begin{center}
\renewcommand{\arraystretch}{1.2}
\begin{tabular}{lccc}
\toprule
$(r,s,t)$ & Operation & Sign $(-1)^{rs+t}$ \\
\midrule
$(0,3,1)$ & $\tilde{m}_2(\tilde{m}_3 \otimes \id)$ & $-1$ \\
$(1,3,0)$ & $\tilde{m}_2(\id \otimes \tilde{m}_3)$ & $-1$ \\
$(0,2,2)$ & $\tilde{m}_3(\tilde{m}_2 \otimes \id \otimes \id)$ & $+1$ \\
$(1,2,1)$ & $\tilde{m}_3(\id \otimes \tilde{m}_2 \otimes \id)$ & $-1$ \\
$(2,2,0)$ & $\tilde{m}_3(\id \otimes \id \otimes \tilde{m}_2)$ & $+1$ \\
\bottomrule
\end{tabular}
\end{center}

These correspond to the five \emph{edges} of the pentagon $K_4$
(not the five vertices, which enumerate the trees for~$m_4$).
The identity is the homotopy-transfer image of \(m_1^2=0\) and the
associativity relation for \(m_2\); it holds for the full transferred
operation, not only for a chosen scalar projection.
\end{computation}

\begin{theorem}[Genus-\texorpdfstring{$0$}{0} transfer tree expansion]
\label{thm:mk-tree-level}
\index{tree-level!Feynman}
For \(k\geq 2\), let \(\mathrm{PBT}_k\) be the set of planar rooted
binary trees with \(k\) leaves.  The operation transferred from a
differential graded algebra \((\mathcal A,m_1,m_2)\) along
\((\pi,\iota,h)\) is
\[
\widetilde m_k
=\sum_{\Gamma\in\mathrm{PBT}_k}
\epsilon_\Gamma\,\pi\,\mu_\Gamma(h,m_2)\,\iota^{\otimes k}.
\]
Here every internal vertex of \(\Gamma\) is labelled by \(m_2\), every
internal edge by \(h\), and \(\epsilon_\Gamma\) is the Koszul sign
from the bar suspension convention.  A tree in \(\mathrm{PBT}_k\) has
\(k-1\) binary vertices, \(k-2\) internal homotopy edges, and
\(\ell(\Gamma)=0\).  Planar rooted trees have automorphism factor \(1\);
after quotienting to non-planar trees, each graph is weighted by
\(|\Aut(\Gamma)|^{-1}\).
\end{theorem}

\begin{proof}
This is the Kadeishvili homotopy transfer formula, equivalently the
sum over vertices of the Stasheff associahedron \(K_k\).  The count of
vertices and internal edges is the elementary count for a rooted
binary tree with \(k\) leaves.  The loop number vanishes because a tree
has \(h^1=0\).
\end{proof}

\begin{theorem}[All-genus stable-graph expansion]
\label{thm:mk-general-structure}
\index{Feynman transform|textbf}
Assume \(\cA\) is equipped with the modular chiral structure of
Theorem~\ref{thm:bar-modular-operad} and with Hodge homotopy
retractions whose edge kernels \(P_g(z,w)\) have diagonal residue
\(1\).  Then the genus-\(g\), \(k\)-input operation has the stable-graph
expansion
\[
m_k^{(g)}(\phi_1,\ldots,\phi_k)
=
\sum_{\Gamma\in \StGraph_{g,k}}
\frac{\epsilon(\Gamma)}{|\Aut(\Gamma)|}
\operatorname{Contr}_\Gamma
\left(
\bigotimes_{v\in V(\Gamma)} C_{g_v,n_v};
\bigotimes_{e\in E(\Gamma)}P_{g(e)}
\right)(\phi_1,\ldots,\phi_k),
\]
where \(\StGraph_{g,k}\) is the set of stable graphs with \(k\) legs
and total genus
\[
g=\sum_{v\in V(\Gamma)}g_v+h^1(\Gamma).
\]
The factor \(1/|\Aut(\Gamma)|\) is the modular-operad automorphism
factor, \(C_{g_v,n_v}\) is the genus-\(g_v\), \(n_v\)-point chiral
operation at the vertex \(v\), and \(\operatorname{Contr}_\Gamma\)
contracts the tensor factors along the propagator kernels on the
internal edges.  The genus-\(0\), \(h^1=0\) specialization is the
transfer-tree expansion of Theorem~\ref{thm:mk-tree-level}.
\end{theorem}

\begin{proof}
Theorem~\ref{thm:prism-higher-genus} and
Theorem~\ref{thm:bar-modular-operad} identify the higher-genus bar
complex as an algebra over the Feynman transform of the commutative
modular operad.  By the definition of the Feynman transform
\textup{(}Definition~\textup{\ref{def:relative-feynman-transform}};
Getzler--Kapranov~\cite{GeK98}\textup{)}, its arity-\(k\), genus-\(g\)
component decomposes as
\[
\bar{B}^{(g)}_k = \bigoplus_{\Gamma \in \StGraph_{g,k}}
\frac{1}{|\mathrm{Aut}(\Gamma)|}\;
\bigotimes_{v \in V(\Gamma)}
\overline{\mathcal{M}}_{g_v, n_v}
\;\otimes\; \bigotimes_{e \in E(\Gamma)} \bC_e
\]
where each vertex~$v$ carries genus~$g_v$ and valence~$n_v$, and each
internal edge~$e$ carries a one-dimensional propagator line.  Evaluating
the vertex tensors on the chiral operations \(C_{g_v,n_v}\) and the
edge tensors on the chosen kernels \(P_g\) gives exactly the displayed
contraction.  The automorphism factor is part of the coinvariant
definition of the modular Feynman transform.
\end{proof}

\begin{remark}
The loop number of a stable graph is \(h^1(\Gamma)\).  The curve genus
is \(\sum_v g_v+h^1(\Gamma)\).  If every vertex has genus \(0\), then
the curve genus equals the graph loop number; otherwise vertex genus
contributes independently.
\end{remark}

\section{Forest comparison obstruction}
\label{sec:bphz-ainfinity}

\begin{definition}[BPHZ--\texorpdfstring{$A_\infty$}{A-infinity} comparison obstruction]
\label{def:bphz-comparison-obstruction}
Let \(\mathfrak F(\Gamma)\) be the forest poset of proper nested
subgraphs of a graph~\(\Gamma\), and let
\(\partial_{\mathrm{St}}\) be the boundary operator on the Stasheff
cell whose faces encode the \(A_\infty\) composites.  A choice of
renormalization scheme and graph assignment gives a linear comparison
map
\[
\mathcal R_n\colon C_\bullet(K_n)\longrightarrow
\bC[\mathfrak F(\Gamma)]
\]
from associahedral boundary chains to forest chains.  The obstruction
to identifying the two recursions at arity \(n\) is
\[
\mathfrak b_n
=\mathcal R_n(\partial_{\mathrm{St}}K_n)
-\partial_{\mathrm{forest}}\mathcal R_n(K_n).
\]
The comparison holds for the chosen renormalization scheme exactly when
\(\mathfrak b_n=0\) for all relevant \(n\).
\end{definition}

\begin{example}[Arity-\texorpdfstring{$3$}{3} obstruction]
At arity \(3\), the uncurved \(A_\infty\) relation is
\begin{align*}
&m_1(m_3(\phi_1 \otimes \phi_2 \otimes \phi_3))
+ m_2(m_2(\phi_1 \otimes \phi_2) \otimes \phi_3)
+ m_2(\phi_1 \otimes m_2(\phi_2 \otimes \phi_3)) \\
&+ m_3(m_1(\phi_1) \otimes \phi_2 \otimes \phi_3) + \cdots = 0.
\end{align*}
The two binary composites are the two codimension-\(1\) faces of
\(\overline{M}_{0,4}\) visible in this bracketing.  Applying
\(\mathcal R_3\) sends them to the corresponding two forest
subtractions.  The class \(\mathfrak b_3\) records exactly the
difference between the boundary-orientation signs
\((-1)^{rs+t}\), the graph automorphism factors, and the chosen
renormalization subtraction signs.
\end{example}

\section{Logarithmic edge forms}

\begin{definition}[Bar graph form]
\label{def:bar-graph-form}
For fields \(\phi_1,\ldots,\phi_k\) and an ordered edge set
\(E(\Gamma)\), the bar graph form is
\[
\omega_\Gamma
=\frac{\epsilon(\Gamma)}{|\Aut(\Gamma)|}
\phi_1(z_1)\otimes\cdots\otimes\phi_k(z_k)
\otimes\bigwedge_{(i,j)\in E(\Gamma)}\eta_{ij}.
\]
On a disk chart, with scalar Cauchy kernel
\(G_0(z_i,z_j)=(z_i-z_j)^{-1}\),
\[
\eta_{ij}=d\log(z_i-z_j)=-\frac{dG_0(z_i,z_j)}{G_0(z_i,z_j)},
\qquad
\operatorname{PRes}_{D_{ij}}\eta_{ij}=1.
\]
Thus the bar edge is the logarithmic derivative of the scalar Cauchy
kernel with residue normalization fixed along the collision divisor.
\end{definition}

\label{sec:summary-feynman-unity}
